\documentclass[a4paper]{article}

\usepackage{graphicx}
\usepackage{amsmath,amssymb}
\usepackage{minted}
\usepackage{multirow}
\usepackage{float}
\usepackage{todonotes}
\usepackage{forest}
\usepackage{xstring}
\usepackage{xcolor}
\usepackage[most]{tcolorbox}
\usepackage{xparse}
\usepackage{natbib}

\usetikzlibrary{graphs,graphdrawing}
\usegdlibrary{layered}

% for external graphviz
\input{/home/donkarlo/Dropbox/repo/nd_language_ai_project/src/nd_language_ai/formal/computer/programming/tex/latex/code_clips/external_graphviz.tex}

% for tags
\input{/home/donkarlo/Dropbox/repo/nd_language_ai_project/src/nd_language_ai/formal/computer/programming/tex/latex/parts/control_sequence/kind/semtag/semtag.tex}

% for links
\input{/home/donkarlo/Dropbox/repo/nd_language_ai_project/src/nd_language_ai/formal/computer/programming/tex/latex/code_clips/seehere.tex}

\title{Receptive field}
\date{}
\author{Mohammad Rahmani}

\begin{document}
    \maketitle
    a portion of sensory space that can elicit neuronal responses when stimulated. The sensory space can be defined in a single dimension (e.g. carbon chain length of an odorant), two dimensions (e.g. skin surface), or multiple dimensions (e.g. space, time[,] and tuning properties of a visual receptive field). The neuronal response can be defined as firing rate (i.e. number of action potentials generated by a neuron) or include also subthreshold activity (i.e. depolarizations and hyperpolarizations in membrane potential that do not generate action potentials) \seeherefooter{https://en.wikipedia.org/wiki/Receptive_field}.
    
    
    \section{Computational theory}
        \seeherefooter{https://en.wikipedia.org/wiki/Receptive_field#Computational_theory_of_visual_receptive_fields}
        
        \seeherefooter{}
        \bibliographystyle{plainnat}
        \bibliography{/home/donkarlo/Dropbox/repo/nd_shared_research/bibliography/bibliography.bib}
\end{document}